\chapter{Konzeptionierung}
\label{cha:concept}

Während Kapitel~\ref{cha:methodik} den Rahmen festlegt, in dem der Vergleich beider Systeme
untersucht wird, beschreibt dieses Kapitel den Untersuchungsgegenstand selbst. Der Reflexionsbot ist
kein vorgefundenes System, sondern wurde für diese Arbeit eigens entwickelt. Seine Gestaltung ist
damit nicht bloß technische Vorarbeit, sondern der eigentliche didaktische Eingriff, dessen Wirkung
anschließend untersucht wird. Sie wird hier so dokumentiert, dass sich die getroffenen
Entscheidungen nachvollziehen und der Bot reproduzieren lässt.

Das Kapitel gliedert sich in vier Abschnitte. Zunächst werden die funktionalen und
nicht-funktionalen Anforderungen aus Zielsetzung und Forschungsfrage abgeleitet
(Abschnitt~\ref{sec:anforderungen}). Darauf aufbauend beschreibt
Abschnitt~\ref{sec:konzept-reflexionsbot} das didaktische Konzept des Bots, insbesondere das Prinzip
der adaptiven Tiefe seiner Rückfragen. Abschnitt~\ref{sec:konzeption-szenarien} stellt die Gestaltung
der beiden Prozessszenarien dar, die für beide Gruppen die identische Aufgabengrundlage bilden.
Abschnitt~\ref{sec:technische-umsetzung} dokumentiert abschließend die technische Umsetzung von der
Auswahl des Frontend-Frameworks bis zur Anbindung des Sprachmodells.

\section{Anforderungen an das System}
\label{sec:anforderungen}

Aus der in Kapitel~\ref{chap:einleitung} formulierten Zielsetzung sowie der in
Abschnitt~\ref{sec:forschungsfrage} formulierten Forschungsfrage lassen sich funktionale und
nicht-funktionale Anforderungen an den zu entwickelnden Reflexionsbot ableiten.

Funktional muss der Bot erstens auf direkte Lösungen verzichten und stattdessen durch Rückfragen
zur eigenen Argumentation der Studierenden anregen, da genau dieser Verzicht den Kern der
Forschungsfrage bildet. Zweitens muss sich die Tiefe dieser Rückfragen adaptiv an die Qualität der
studentischen Argumentation anpassen, damit weder oberflächliche Antworten unkommentiert bleiben
noch bereits durchdachte Antworten unnötig verzögert werden. Drittens muss der Bot beide
Prozessszenarien (SmartDose, LearnLoop, vgl. Abschnitt~\ref{sec:konzeption-szenarien}) kennen und
inhaltlich referenzieren können, ohne sie beim Namen zu nennen, damit Studierende nicht anhand der
Szenario-Bezeichnung nachschlagen, statt selbst zu argumentieren. Viertens müssen Studierende
mehrere Chat-Sitzungen parallel führen und wechseln können, und fünftens müssen die
Gesprächsverläufe für die spätere qualitative Auswertung exportierbar sein.

Nicht-funktional muss der Bot ohne lokale Installation zugänglich sein, damit er im Rahmen einer
regulären Übung ohne technische Vorbereitung genutzt werden kann. Die Nutzung muss pseudonym erfolgen,
sodass kein Klarname, sondern nur ein selbst gewählter Personal Code erforderlich ist (vgl.
Abschnitt~\ref{sec:datenerhebung}). Bei der Modellwahl war zudem auf ein angemessenes Verhältnis
von Antwortqualität und Kosten zu achten, da mehrere Studierende die Anwendung gleichzeitig während
der Übung nutzen sollten. Schließlich sollten Antworten möglichst unmittelbar sichtbar werden, um
den Gesprächsfluss nicht durch Wartezeiten zu unterbrechen.

Tabelle~\ref{tab:anforderungen} fasst die Anforderungen zusammen und weist jeder Anforderung aus,
worauf sie zurückgeht. Die Zuordnung macht sichtbar, dass die funktionalen Anforderungen FA1 bis FA3
unmittelbar aus den Teilfragen folgen und damit die Untersuchungsbedingung selbst herstellen, während
FA4 und FA5 sowie die nicht-funktionalen Anforderungen den Erhebungs- und Durchführungsrahmen
absichern. Diese Unterscheidung ist für die Bewertung der Umsetzung wesentlich: Ein Fehler in FA1
oder FA2 gefährdet die Gültigkeit des Vergleichs, ein Fehler in NFA3 lediglich den Komfort.

\begin{table}[!t]
\centering
\caption{Funktionale und nicht-funktionale Anforderungen an den Reflexionsbot}
\label{tab:anforderungen}
\renewcommand{\arraystretch}{1.3}
\setlength{\tabcolsep}{8pt}
\begin{tabular}{>{\raggedright\arraybackslash}p{0.07\textwidth}
                >{\raggedright\arraybackslash}p{0.43\textwidth}
                >{\raggedright\arraybackslash}p{0.30\textwidth}}
\toprule
\textbf{Nr.} & \textbf{Anforderung} & \textbf{Begründet durch} \\
\midrule
\multicolumn{3}{@{}l}{\textit{Funktionale Anforderungen}} \\
\midrule
FA1 & Verzicht auf direkte Lösungen, stattdessen Rückfragen zur eigenen Argumentation &
TF1, Kern der Vergleichs\-bedingung \\
\midrule
FA2 & Adaptive Tiefe der Rückfragen, abhängig von der Qualität der studentischen Antwort &
TF2, Scaffolding-Prinzip \\
\midrule
FA3 & Kenntnis beider Szenarien, ohne sie beim Namen zu nennen &
TF2, verhindert Nachschlagen statt Argumentieren \\
\midrule
FA4 & Mehrere parallele Chat-Sitzungen mit Wechselmöglichkeit &
Durchführung, zwei Szenarien je Person \\
\midrule
FA5 & Exportierbarkeit der Gesprächsverläufe &
Datengrundlage der qualitativen Kodierung \\
\midrule
\multicolumn{3}{@{}l}{\textit{Nicht-funktionale Anforderungen}} \\
\midrule
NFA1 & Zugänglich ohne lokale Installation &
Einsatz in regulärer Übung ohne Vorbereitung \\
\midrule
NFA2 & Pseudonyme Nutzung über selbst gewählten Personal Code &
Forschungsethik, Prä-Post-Verknüpfung \\
\midrule
NFA3 & Angemessenes Verhältnis von Antwortqualität und Kosten &
Gleichzeitige Nutzung durch mehrere Studierende \\
\midrule
NFA4 & Unmittelbar sichtbare Antworten &
Erhalt des Gesprächsflusses \\
\bottomrule
\end{tabular}
\end{table}

\section{Konzept des Reflexionsbots}
\label{sec:konzept-reflexionsbot}

Der entwickelte Reflexionsbot ist kein generisches Frage-Antwort-System. Er ist ein
spezialisierter Reflexionscoach, der Studierende bei der Bearbeitung von zwei
Szenarien zur Prozessmodellierung aus dem Software Engineering begleitet. Die
Szenarien bilden realistische Entscheidungssituationen ab, etwa Rollen-, Daten- und
Zeitplanentscheidungen in einem regulierten ML-Projekt sowie Architektur- und
Kommunikationsentscheidungen in einem verteilten Team, und dienen als fester
Kontext, auf den sich alle Rückfragen des Bots beziehen (vollständige
Szenariobeschreibungen siehe \hyperref[app:szenarien]{Anhang~A.4}).

Im Zentrum des Konzepts steht das in Abschnitt~\ref{sec:relevante-technologien-und-konzepte}
eingeführte Prinzip von Scaffolding und sokratischem Fragen: Der Bot liefert keine
fertigen Lösungen, sondern stellt gezielte Rückfragen, die Studierende zur Begründung
und Reflexion ihrer eigenen Prozessentscheidungen anregen. Damit unterscheidet er sich
bewusst von einem generischen \acrshort{ki}-System, das eine Frage typischerweise direkt
beantworten würde. Genau dieser Unterschied bildet den Kern der Forschungsfrage dieser
Arbeit.

Ein zentrales Gestaltungsmerkmal ist die \emph{adaptive Tiefe} der Bot-Antworten. Die
Ausführlichkeit und Qualität der Rückmeldung skaliert mit der Qualität der studentischen
Argumentation, wie Abbildung~\ref{fig:adaptive-tiefe} zusammenfasst.

\begin{itemize}
    \item \textbf{Vage Antwort} (z.\,B. Nennung einer Option ohne Begründung): Der Bot
        stellt genau eine kurze, gezielte Rückfrage, um die Begründung herauszuarbeiten.
    \item \textbf{Teilweise Begründung} (Studierende nennen eine Entscheidung mit einer
        Begründung, lassen aber Risiken oder Randbedingungen aus): Der Bot bestätigt kurz
        den richtigen Teil und stellt ein bis zwei fokussierte Rückfragen zu den fehlenden
        Aspekten.
    \item \textbf{Starke Argumentation} (Abwägungen, Randbedingungen, Konsequenzen werden
        benannt): Der Bot bestätigt den stärksten Punkt explizit, ergänzt eine Perspektive,
        die nicht genannt wurde, und stellt eine weiterführende Frage.
    \item \textbf{Herausragende Antwort} (mehrperspektivische Abwägung, konkrete
        Randbedingungen, Folgeeffekte werden antizipiert): Der Bot begegnet den
        Studierenden auf Augenhöhe, benennt das stärkste Gegenargument und stellt eine
        Frage, die die Kernannahme herausfordert.
\end{itemize}

\begin{figure}[htb]
\centering
\begin{tikzpicture}[
	start/.style={rectangle, draw=black!55, fill=black!7, rounded corners, align=center,
	              text width=5.4cm, minimum height=0.8cm, font=\small},
	stufe/.style={rectangle, draw=codeblue, rounded corners, align=left, anchor=north,
	              text width=3.05cm, minimum height=3.1cm, font=\scriptsize, inner sep=5pt},
	arr/.style={-{Latex[length=1.8mm]}, draw=black!65}
	]
\node[start] (ant) at (0,0) {Studentische Antwort};

\node[stufe, fill=codeblue!6]  (s1) at (-5.55,-1.45)
  {\textbf{Vage Antwort}\\[1pt] Option ohne Begründung\\[5pt]
   $\rightarrow$ genau \emph{eine} kurze, gezielte Rückfrage};
\node[stufe, fill=codeblue!16] (s2) at (-1.85,-1.45)
  {\textbf{Teilweise Begründung}\\[1pt] Risiken und Randbedingungen fehlen\\[5pt]
   $\rightarrow$ kurze Bestätigung und ein bis zwei fokussierte Rückfragen};
\node[stufe, fill=codeblue!28] (s3) at (1.85,-1.45)
  {\textbf{Starke Argumentation}\\[1pt] Abwägungen und Konsequenzen benannt\\[5pt]
   $\rightarrow$ stärksten Punkt bestätigen, Perspektive ergänzen, weiterführende Frage};
\node[stufe, fill=codeblue!42] (s4) at (5.55,-1.45)
  {\textbf{Herausragende Antwort}\\[1pt] mehrperspektivisch, Folgeeffekte antizipiert\\[5pt]
   $\rightarrow$ auf Augenhöhe, stärkstes Gegenargument, Kernannahme hinterfragen};

\foreach \n in {s1,s2,s3,s4} \draw[arr] (ant.south) -- (\n.north);
\end{tikzpicture}
\caption[Adaptive Tiefe der Bot-Antworten]{Adaptive Tiefe der Bot-Antworten. Mit steigender Qualität
der studentischen Argumentation nimmt die Ausführlichkeit der Rückmeldung zu (zunehmend dunklere
Färbung), während bei einer vagen Antwort bewusst nur eine einzelne Rückfrage erfolgt.}
\label{fig:adaptive-tiefe}
\end{figure}

Diese gestufte Unterstützung setzt das Scaffolding-Prinzip technisch um: Hilfestellung
wird nicht pauschal gewährt, sondern in Abhängigkeit vom gezeigten Kompetenzniveau
zurückgenommen, sobald Studierende eigenständiger argumentieren. Zusätzlich fokussiert
der Bot pro Gesprächsrunde jeweils eine logische Schwachstelle (z.\,B. zunächst Risiko,
danach Rollen und Stakeholder, dann Abhängigkeiten, dann Zeitplanrealismus), sodass sich
die Reflexion schrittweise vertieft, statt alle Schwächen gleichzeitig zu adressieren.
Um zu verhindern, dass Studierende Antworten anhand der Szenario-Bezeichnung erraten
oder nachschlagen, statt selbst zu argumentieren, nennt der Bot die Szenarien in seinen
Antworten nie explizit beim Namen.

Schließlich antwortet der Bot durchgängig auf Englisch, unabhängig davon, in welcher Sprache die
Studierenden schreiben. Hintergrund ist, dass die Lehrveranstaltung \glqq Grundlagen des Software
Engineering 1\grqq{} auch internationale Studierende umfasste. Eine einheitlich englische
Kommunikation stellte sicher, dass auch diese Teilnehmenden eingebunden waren und die Rückfragen des
Bots verstehen konnten. Aus demselben Grund wurden die Aufgabenblätter (\hyperref[app:szenarien]{Anhang~A.4}) und die
Fragebögen (\hyperref[app:fragebogen]{Anhang~A.5}) auf Englisch ausgegeben.

Abbildung~\ref{fig:reflexionsbot-ui} zeigt die Oberfläche des Reflexionsbots anhand eines
beispielhaften Gesprächsverlaufs zur Datensatzentscheidung aus dem Szenario SmartDose. Über die
Seitenleiste lassen sich mehrere Chats parallel anlegen und der Gesprächsverlauf als Textdatei
herunterladen, während das Eingabefeld im Hauptbereich für die Kommunikation mit dem Bot genutzt
wird. Der Verlauf veranschaulicht zugleich die adaptive Tiefe. Auf die einleitende Frage nach der zu
wählenden Option antwortet der Bot nicht mit einer Empfehlung, sondern mit einer einzelnen gezielten
Rückfrage nach den Risiken. Auf die knappe Antwort \glqq Maybe we could get into legal trouble\grqq{}
folgt eine kurze Bestätigung und eine fokussierte Nachfrage zur nicht gewählten Option. Erst als die
studentische Antwort Zulassungsfähigkeit und Modellqualität explizit gegeneinander abwägt, bestätigt
der Bot den stärksten Punkt ausdrücklich und ergänzt eine weiterführende Frage.

Im Vergleich zu bestehenden textgenerativen Reflexionswerkzeugen wie \glqq Reflect Me\grqq{}
\cite{mollerReflectMeTextgenerative2025}, das allgemeine Reflexionsimpulse zu frei wählbaren
Themen gibt, ist der hier entwickelte Reflexionsbot bewusst enger gefasst: Er kennt die konkreten
Entscheidungssituationen beider Szenarien im Detail und kann seine Rückfragen daher gezielt auf die
jeweilige Prozessentscheidung sowie die dort relevanten Randbedingungen zuschneiden, statt allgemein
gehaltene Reflexionsfragen zu stellen. Diese Szenario-Spezifität ist die Voraussetzung dafür, dass
die adaptive Tiefe (s.\,o.) inhaltlich präzise auf einzelne Argumentationslücken reagieren kann.

\begin{figure}[htb]
\centering
\includegraphics[width=\textwidth]{BotImage.png}
\caption[Oberfläche des Reflexionsbots mit Beispielverlauf]{Oberfläche des Reflexionsbots
(Streamlit-Anwendung) mit einem beispielhaften Gesprächsverlauf zur Datensatzentscheidung aus dem
Szenario SmartDose. Der Verlauf wurde zur Illustration der adaptiven Tiefe nachgestellt und stammt
nicht aus dem Mini-Experiment.}
\label{fig:reflexionsbot-ui}
\end{figure}


\section{Konzeption der Szenarien}
\label{sec:konzeption-szenarien}

Neben dem Reflexionsbot selbst war die Gestaltung der beiden Prozessszenarien SmartDose und
LearnLoop ein eigenständiger Teil der Konzeption, da die Qualität der möglichen Reflexion maßgeblich
von der Aufgabenstellung abhängt, auf die sich diese Reflexion bezieht (vollständige
Szenariobeschreibungen siehe \hyperref[app:szenarien]{Anhang~A.4}).

Beiden Szenarien liegt dasselbe Gestaltungsprinzip zugrunde: Sie bestehen aus jeweils drei
Dilemma-Situationen, die bewusst keine eindeutig richtige Lösung besitzen, sondern eine Abwägung
zwischen mehreren vertretbaren, aber jeweils mit Nachteilen behafteten Optionen verlangen. Die
Aufgabenstellung schließt Ausweichantworten wie \glqq we would do both\grqq{} ausdrücklich aus, um
eine tatsächliche Entscheidung zu erzwingen. Um diese Abwägung inhaltlich zu verankern, sind beide
Szenarien um reale regulatorische und organisationale Zwänge herum konstruiert: SmartDose stellt
Patientensicherheit und Entwicklungsgeschwindigkeit gegen die Vorgaben der EU-Medizinprodukteverordnung
(MDR), LearnLoop stellt Projektfortschritt in einem verteilten Team gegen Datenschutzanforderungen
(FERPA) und massive Zeitzonenunterschiede. Beide Konflikte sind so gewählt, dass eine rein technische
Betrachtung nicht ausreicht, sondern rechtliche, ethische und organisationale Gesichtspunkte
gegeneinander abgewogen werden müssen.

Jedes Szenario folgt zudem demselben dreistufigen Aufgabenaufbau, den
Abbildung~\ref{fig:aufgabenaufbau} zusammenfasst. Zunächst analysieren Studierende
in Task~1, welche Scrum-Elemente im veränderten Kontext noch funktionieren (Verständnis der
Ausgangslage), anschließend treffen und begründen sie in Task~2 Entscheidungen zu den drei
Dilemma-Situationen (eigentliche Entscheidungsfindung), und schließlich reflektieren sie in Task~3,
welche Entscheidung am schwersten fiel und warum (Metareflexion). Bei Team AI folgt darauf ein
eigenständiger Gegencheck mit einem generischen KI-System, beim Reflexionsbot erfolgt dieser
Gegencheck bereits im laufenden Dialog mit dem Bot selbst.

\begin{figure}[htb]
\centering
\begin{tikzpicture}[
	task/.style={rectangle, draw=black!55, fill=black!7, rounded corners, align=center,
	             text width=7.2cm, minimum height=0.95cm, font=\small, inner sep=5pt},
	ai/.style={rectangle, draw=codegreen, line width=0.8pt, fill=codegreen!12, rounded corners,
	           align=center, anchor=north, text width=4.9cm, minimum height=1.5cm,
	           font=\small, inner sep=5pt},
	rb/.style={rectangle, draw=codeblue, line width=0.8pt, fill=codeblue!12, rounded corners,
	           align=center, anchor=north, text width=4.9cm, minimum height=1.5cm,
	           font=\small, inner sep=5pt},
	arr/.style={-{Latex[length=2mm]}, draw=black!65}
	]
\node[task] (t1) at (0,0)     {\textbf{Task 1 -- Analyse}\\ Welche Scrum-Elemente funktionieren im veränderten Kontext?};
\node[task] (t2) at (0,-1.85) {\textbf{Task 2 -- Decide}\\ Entscheidung und Begründung zu drei Dilemma-Situationen};
\node[task] (t3) at (0,-3.7)  {\textbf{Task 3 -- Reflect}\\ Welche Entscheidung fiel am schwersten und warum?};

\node[ai] (ai) at (-2.85,-5.5) {\textbf{Team AI}\\ eigenständiger Gegencheck mit einem generischen KI-System};
\node[rb] (rb) at (2.85,-5.5)  {\textbf{Reflexionsbot}\\ Gegencheck bereits im laufenden Dialog mit dem Bot};

\draw[arr] (t1) -- (t2);
\draw[arr] (t2) -- (t3);
\draw[arr] (t3.south) -- (ai.north);
\draw[arr] (t3.south) -- (rb.north);
\end{tikzpicture}
\caption[Dreistufiger Aufgabenaufbau der Szenarien]{Dreistufiger Aufgabenaufbau der Szenarien. Beide
Gruppen bearbeiteten Task~1 bis Task~3 mit identischen Aufgabenblättern. Der einzige Unterschied lag
in der abschließenden Gegenargumentation.}
\label{fig:aufgabenaufbau}
\end{figure} Dieser dreistufige Aufbau war notwendig,
um die in Abschnitt~\ref{sec:datenanalyse} beschriebene Kodierung nach Reflexionsstufen überhaupt an
einer vergleichbaren Aufgabenstruktur vornehmen zu können.

\section{Technische Umsetzung}
\label{sec:technische-umsetzung}

Abbildung~\ref{fig:architektur} gibt einen Überblick über die technische Architektur, bevor die
einzelnen Entwurfsentscheidungen im Detail erläutert werden: Studierende interagieren über den
Browser mit der Streamlit-Anwendung, die Eingaben entgegennimmt, den Sitzungszustand verwaltet und
bei jeder Anfrage die Systemnachricht zusammen mit dem bisherigen Gesprächsverlauf an die
Chat-Completions-API von OpenAI übergibt. Die Anwendung selbst läuft containerisiert in Docker.

\begin{figure}[htbp]
\centering
\begin{tikzpicture}[
	box/.style={rectangle, draw, rounded corners, align=center, minimum height=1.1cm,
	            text width=4.8cm, font=\small, inner sep=6pt},
	extern/.style={box, draw=black!55, fill=black!7},
	eigen/.style={box, draw=codeblue, line width=0.9pt, fill=codeblue!12},
	arr/.style={{Latex[length=2mm]}-{Latex[length=2mm]}, draw=black!70},
	lbl/.style={font=\scriptsize, fill=white, inner sep=2.5pt, align=center}
	]
	\node[extern] (user) {Studierende (Browser)};
	\node[eigen, below=2.9cm of user] (streamlit) {Streamlit-App\\(Frontend und \texttt{session\_state})};
	\node[extern, below=2.9cm of streamlit] (api) {OpenAI Chat-Completions-API\\(\texttt{gpt-4o-mini})};

	\draw[arr] (user) -- node[lbl] {Eingabe / Anzeige} (streamlit);
	\draw[arr] (streamlit) -- node[lbl] {Systemprompt + Verlauf /\\Streaming-Antwort} (api);

	\begin{scope}[on background layer]
	\node[draw=codeblue!55, dashed, rounded corners, fill=codeblue!4,
	      fit=(streamlit), inner sep=13pt,
	      label={[font=\scriptsize, text=codeblue!80]above right:Docker-Container}] {};
	\end{scope}
\end{tikzpicture}
\caption{Vereinfachte Architektur des Reflexionsbots}
\label{fig:architektur}
\end{figure}

\subsection{Frontend-Auswahl: Streamlit vs. Gradio}
Für die Umsetzung des Chat-Interfaces standen Streamlit \cite{streamlitSessionState2024} und
Gradio \cite{gradioChatInterface2024} zur Auswahl, da beide Frameworks eine schnelle Entwicklung
von Python-basierten Web-Anwendungen ohne separaten Frontend-Stack ermöglichen.
Tabelle~\ref{tab:streamlit-gradio} stellt beide Frameworks anhand der für den Reflexionsbot
relevanten Kriterien gegenüber \cite{streamlitSessionState2024, gradioChatInterface2024}.

\begin{table}[htb]
\centering
\caption{Vergleich von Streamlit und Gradio als Frontend-Framework}
\label{tab:streamlit-gradio}
\renewcommand{\arraystretch}{1.3}
\setlength{\tabcolsep}{8pt}
\begin{tabular}{>{\raggedright\arraybackslash}p{0.31\textwidth}
                >{\raggedright\arraybackslash}p{0.29\textwidth}
                >{\raggedright\arraybackslash}p{0.25\textwidth}}
\toprule
\textbf{Kriterium} & \textbf{Streamlit} & \textbf{Gradio} \\
\midrule
Vorgefertigte Chat-Komponenten &
Nicht enthalten, Chat-Oberfläche muss selbst aufgebaut werden &
Vorhanden, dadurch schneller einsatzbereiter Prototyp \\
\midrule
Mehrere parallele Chat-Sitzungen &
Über \texttt{session\_state} flexibel umsetzbar &
Nicht für diesen Anwendungsfall vorgesehen \\
\midrule
Export der Gesprächsverläufe &
Direkter Zugriff auf den Sitzungszustand, frei exportierbar &
Nur eingeschränkt vorgesehen \\
\midrule
Kontrolle über Layout und Zustandsverwaltung &
Umfangreich, da als vollwertige Webanwendung strukturierbar &
Eingeschränkt, komponentenbasiert \\
\midrule
Einbettung in ein größeres Analyse-Dashboard &
Gut geeignet &
Primär auf einzelne Modell-Demos zugeschnitten \\
\midrule
Entwicklungs\-geschwindig\-keit für einfachen Prototyp &
Etwas höherer Aufwand &
Sehr schnell \\
\bottomrule
\end{tabular}
\end{table}

Gradio ist stärker auf einzelne Modell-Demos zugeschnitten und bietet vorgefertigte
Chat-Komponenten, die für einen einfachen Chatbot-Prototyp schneller einsatzbereit sind. Da der
Reflexionsbot jedoch mehrere parallele Chat-Sitzungen verwalten, Gesprächsverläufe für die
spätere Auswertung exportieren und sich potenziell in ein größeres Analyse-Dashboard einbetten
lassen sollte, fiel die Wahl auf Streamlit, das bei diesen Anforderungen mehr Kontrolle über
Layout und Zustandsverwaltung bietet.

\subsection{Implementierung}
Die Anwendung wurde vollständig in Python umgesetzt. Der Chatverlauf jeder Sitzung wird im
Streamlit-\texttt{session\_state} als Liste von Nachrichten verwaltet. Über eine Seitenleiste
lassen sich mehrere Chats parallel anlegen, wechseln und als Textdatei herunterladen. Der
Titel eines Chats wird automatisch aus der ersten Eingabe der Studierenden abgeleitet.

Die Entscheidung für die OpenAI-\acrshort{api} gegenüber einem lokal betriebenen Modell begründet
sich in mehreren Faktoren: Der Betrieb leistungsfähiger \glspl{llm} erfordert erhebliche
Rechenressourcen, die im Rahmen dieser Arbeit nicht verfügbar waren. Zusätzlich bieten Modelle der
\acrshort{gpt}-4-Familie eine hohe Qualität bei Sprachverarbeitung und Gesprächsführung
\cite{openaiGPT4TechnicalReport2023}, und die REST-basierte Schnittstelle lässt sich ohne aufwendige
Infrastruktur direkt in Webanwendungen integrieren \cite{openaiChatCompletionsAPI2024}.

Die Anbindung an das Sprachmodell erfolgt über die in
Abschnitt~\ref{sec:relevante-technologien-und-konzepte} beschriebene
Chat-Completions-\acrshort{api} von OpenAI \cite{openaiChatCompletionsAPI2024}, wobei als
Modell \texttt{gpt-4o-mini} verwendet wird. Die Wahl der ressourcenschonenden mini-Variante
gegenüber dem größeren \texttt{gpt-4o} folgt der in Abschnitt~\ref{sec:anforderungen} formulierten
Anforderung eines angemessenen Verhältnisses von Antwortqualität und Kosten, da während der Übung
mehrere Studierende die Anwendung gleichzeitig nutzen sollten. Für die vom Bot benötigten Rückfragen
war die Qualität der mini-Variante ausreichend. Bei jeder Anfrage wird die
Systemnachricht (vgl. Abschnitt~\ref{sec:konzept-reflexionsbot}; vollständiger Wortlaut in
\hyperref[app:systemprompt]{Anhang~A.7}) zusammen mit dem vollständigen bisherigen Gesprächsverlauf an das Modell übergeben, damit der
Kontext über mehrere Gesprächsrunden erhalten bleibt. Die Antwort wird gestreamt
(\texttt{stream=True}) an die Oberfläche ausgegeben, sodass Studierende die Antwort des
Bots bereits während der Generierung lesen können, statt auf die vollständige Antwort
warten zu müssen. Abbildung~\ref{fig:api-call} zeigt diesen zentralen Mechanismus anhand des
tatsächlich eingesetzten Codes.

\begin{figure}[htbp]
\centering
\begin{lstlisting}[language=Python]
# The Chat Completions API is stateless: it has no memory of earlier
# calls, so the system prompt and the full message history are sent
# again with every request to keep the model's context consistent.
conversation = [{"role": "system", "content": SYSTEM_PROMPT}]
conversation.extend(
    {"role": m["role"], "content": m["content"]}
    for m in current_messages
)
stream = client.chat.completions.create(
    model=st.session_state["openai_model"],
    messages=conversation,
    stream=True,
)
response = st.write_stream(stream)
\end{lstlisting}
\caption{Aufbau der Nachrichtenliste und Streaming-Aufruf der Chat-Completions-API}
\label{fig:api-call}
\end{figure}

Für die Durchführung des Mini-Experiments wurde die Anwendung containerisiert
(Docker) und öffentlich erreichbar bereitgestellt, sodass Studierende unabhängig von
lokalen Python-Installationen darauf zugreifen konnten.
